\section{The experimental problem}\label{sec:problem}
A causal experiment does not present an observer with a parameter manifold. It presents preparations, admissible interventions and reports. Geometry becomes relevant only after specifying which future reports are to be predicted and which information the observer has actually acquired. There is a further obstruction: a small terminal codebook need not be recursively selectable after earlier reports have been discarded. The mother problem considered here is to connect these two facts to executable comparison of experiments with an explicit resource account.

The central result, Theorem~\ref{thm:transfer}, has two sides. A static cover preserving a budget-independent relation, together with physical-law block stability and robust candidate moments, constructs a finite-state predictor rounded after every report. In the other direction, a dominated continuation cut turns the conditional responses to future reports into a random element of a separable function space. Its actual acquired law supplies a lower bound on every retained alphabet crossing that cut. Terminal quantization is the last member of this family of bounds, not a replacement for it. The distinction produces an explicit online/checkpoint separation, using a classical random-access argument.

A separate difficulty concerns acquisition itself. A geometric lower bound for one fixed exploration does not determine the value of optimizing the exploration. Theorem~\ref{thm:adaptive} addresses this optimization for a class of progressive binary sensors. A single greedy scale sequence controls both the optimal acquisition allocation and the physical-law small-ball bound. Consequently adaptive sensing, online compression and terminal checkpoint compression have the same order under joint acquisition and retained-state constraints. Ratios may vary arbitrarily with scale and coordinate; some factors may be Lebesgue intervals while others are singular. Theorem~\ref{thm:joint} places acquisition, memory, calibration and an attained simulation defect in one common task. The second realization, Theorem~\ref{thm:brier}, concerns squared filtering loss for a primitive sparse hidden-state model with informative missingness and Gaussian reports.

These conclusions have different logical statuses. The block construction extends earlier sufficient results. Continuation quantization is a general obstruction with a product-domination hypothesis, not a sufficiency theorem for recursive realizability. The adaptive theorem optimizes the sensing class below, not arbitrary interventions. Quantization, conditional projection, random-access entropy and comparison inequalities are credited as antecedents. The technical companion~\cite{companion} retains the singular nonreset, stopped, calibration, output-grid and interface results with their original hypotheses.

\subsection{Kernels, interfaces and resources}
All measurable spaces are standard Borel. An experiment is specified by a preparation $\mu_\lambda$, hidden states $X_t$, actions $A_t$, reports $Y_{t+1}$, and jointly measurable probability kernels
\begin{equation}\label{eq:kernel}
 K_t^\lambda(dx',dy,dc\mid x,a).
\end{equation}
Here $c$ records nonnegative raw-call, time and other declared resource increments. Positivity means nonnegativity and normalization. Dirac operations, failure symbols and missing reports are allowed. A visible history $h_t$ consists of the preparation report and the intervening actions, reports and public costs. Legal action sets have a declared measurable graph. Strategies are parameter-independent Borel kernels on those sets; stopping rules use the visible filtration. They cannot consult future reports or the unknown $\lambda$.

A continuation interface is a measurable function $v_t(h_t)$ which specifies legal future actions, available resources and any publicly supplied phase. A predictive realization is a standard Borel state $S_t=s_t(h_t)$ determining that interface, every declared future-test expectation, and jointly Borel report and update kernels $J_t,F_t$. When a phase is kept separate, write the state space as a measurable disjoint union $S=\coprod_{v\in V}S_v$ over a \emph{finite} public phase set. All comparison pairs at a given time lie in the same fiber $S_v$. Updates map both members into the same next fiber, whose phase is a declared function of the old phase, action and common report. No unretained data-dependent phase is public for free. An unbounded or data-dependent external interface requires a separately stated conditional theorem and is not covered by the alphabet count below.

A predictor retains at most $M$ data labels at each declared cut. Its update uses only its current retained tuple, current action/report and allowed fresh randomness. Any controller, simulator or clock register consulted later is counted separately or included in the label tuple. Temporary workspace must be erased before the next cut; program constants must be known, parameter-independent data with a charged description. The notation $M$ is an alphabet cardinality, not total arithmetic space. If a public finite phase has $Q$ states, the full serial cardinality is at most $MQ$; a lower bound at a cut uses the \emph{entire} accessible data-dependent tuple. Deterministic time is fixed in the main prediction statements and is not a hidden observation channel.

For the general transfer theorem fix a preparation and a legal exploration independently of the prediction encoder. Algorithm coins are independent of that experiment. In the adaptive acquisition theorem the exploration is instead optimized together with the encoder; its lower bounds are proved anew and do not condition an encoder-dependent law as though it were fixed.

\subsection{Executable tests and a measurable predictive quotient}
An executable test is a legal finite continuation followed by a bounded function of its scored reports. Two histories are equivalent if their continuation interfaces and all test expectations agree. A set-theoretic quotient need not be a suitable measurable state. We record a sufficient realization result.

\begin{proposition}[Continuous-instrument realization]\label{prop:quotient}
Let $\mathcal B$ be compact metrizable and let $\mathcal V\subset C(\mathcal B)$ be the separable closed linear span of constants and the continuous evaluations of a determining family of executable tests and interface coordinates. All kernel, density, update and instrument versions below are jointly Borel also in the action variable on the declared standard Borel action--report domain. For each fixed action, let its report space be Polish, with reference measure $m_a$, and let $g_a(\pi,y)$ be a jointly continuous density. On the open set $D_a=\{(\pi,y):g_a(\pi,y)>0\}$, with the relative product topology, suppose the normalized update $B_a$ is continuous. For every $f\in\mathcal V$, require the instrument evaluation $g_a(\pi,y)f(B_a(\pi,y))$, with its specified value at density zero, to belong to $\mathcal V$ as a function of $\pi$ and to be jointly continuous. Then the countable evaluation image is a compact metrizable predictive quotient. Report laws and updates descend, and the latter are continuous on the descended open positive-density domain. Every standard Borel statistic measurably determining all these evaluations measurably determines the quotient.
\end{proposition}
\begin{proof}
Normalize a countable dense family of evaluations and map $\mathcal B$ into $[-1,1]^{\N}$. The image $S$ is compact, hence Borel. Equality of its coordinates gives equality on $\mathcal V$ by uniform approximation and hence on the determining tests. Conversely, equality of all future tests and interface coordinates gives equality on their closed linear span. Thus the fibers are exactly predictive equivalence classes, not a finer partition carrying unrelated coordinates. Taking $f=1$ shows that density values agree on each fiber. Dividing equal instrument values by that positive common density shows that the evaluation of the updated belief agrees on the fiber.

The evaluation map is a continuous surjection from a compact space to a Hausdorff space, hence is perfect. Its product with the report identity is also a quotient map: closedness follows by taking a convergent subnet in the compact belief domain whenever an image net converges. The instrument coordinates and their positive-density ratios therefore descend continuously. The positive-density domain is open in $S\times Y$; no update is defined by division at a null report. For measurability jointly across actions, the evaluation map admits a Borel section. Indeed choose nested finite closed covers of the compact domain with diameters tending to zero and, at each stage, the first member meeting the fiber and the previously selected members. Each such test is Borel because the image of the resulting fixed compact intersection is closed. The selected nested intersections have a unique limit in the fiber, and their shrinking diameters make that limit a Borel function of the image point. Composing the assumed joint versions with this section supplies joint Borel versions of all descended kernels. A common jointly Borel null-report extension, when needed on wrong candidates, is additional data.

For minimality, collect the measurable coordinate factors through the given statistic into a measurable product map. It takes values in the Borel set $S$ on the domain in question; extend outside that set by a fixed point. Countably many almost-sure factors can be imposed on a single common full-measure history set. This proves measurable, rather than merely set-theoretic, minimality.
\end{proof}

Applications may verify a predictive realization directly. The metric used for a quantitative estimate need only induce its standard Borel structure; it need not be the compact topology in Proposition~\ref{prop:quotient}.

\subsection{One prediction task}
Let $Y$ be a bounded terminal vector in a finite-dimensional Euclidean space with a specified weighted norm. A finite independent terminal probe index is interpreted as the vector of its separately executable conditional means, with the probe probabilities as weights; it is not a simultaneous counterfactual measurement. At the decision cut $T$ put
\begin{equation}\label{eq:task}
 P_T=\E[Y\mid\mathcal H_T],\qquad B_T=\E\norm{Y-P_T}^2,
 \qquad \nu_T=\law(P_T).
\end{equation}
For a finite Borel measure $\nu$ on a separable Hilbert space define
\begin{equation}\label{eq:quantization}
 q_M(\nu)=\inf_{c_1,\ldots,c_M}\int\min_i\norm{x-c_i}^2\,\nu(dx).
\end{equation}
The measure may be atomic, singular or a subprobability. A checkpoint algorithm inspects $\mathcal H_T$ once before retaining $M$ labels. An online algorithm retains at most $M$ labels at every earlier cut and cannot reread a report. In both cases the terminal outcome is revealed only after the prediction. Denote the optimal risks by $R_\cp(T,M)$ and $R_\on(T,M)$.
